\documentclass[11pt]{article}
\usepackage[margin=1in]{geometry}
\usepackage[T1]{fontenc}
\usepackage[hidelinks]{hyperref}
\setlength{\parindent}{1.5em}\setlength{\parskip}{6pt}\pagestyle{empty}
\begin{document}

\noindent To the Editors\\
\emph{Natural Language Processing Journal}

\vspace{6pt}
\noindent\textbf{Re: Original research submission --- ``Evaluating LLM-Based Credit Rating Systems: A Pre-Registered Specification-Curve Analysis.''}

\vspace{6pt}
\noindent Dear Editors,

We are pleased to submit the above manuscript for consideration in \emph{Natural Language Processing Journal}. The paper contributes an elicitation-reliability evaluation framework for evaluating the model risk of using large language models as credit-rating systems. Using a pre-registered specification-curve design, we measure whether an LLM-generated credit rating is stable when the underlying financial information is held fixed but reasonable elicitation choices vary. Across a fixed 90-item battery crossed with a balanced grid of model version, temperature, prompt paraphrase, output format, few-shot exemplars, and presentation across two providers ($8{,}640$ elicitations, four model snapshots), no controllable elicitation axis explains more than about $5\%$ of rating variance, pure run-to-run seed variation is only about $8\%$, the instability persists on the temperature-0-honored subgrid (rating-invariance permutation $p=0.0005$), and $25.2\%$ of matched elicitation pairs disagree at the investment-grade/high-yield boundary. The pattern persists on 31 perturbed and fingerprint-screened real issuers under a pre-registered fingerprint-identification gate, and a post-hoc reasoning-architecture arm (DeepSeek-R1) shows comparable instability.

The work fits the journal's scope---intelligent systems applied to business and finance---and speaks directly to the deployment of LLMs in financial decision support. Beyond the diagnosis, we propose a candidate operational control: a \emph{specification-instability score}, computed from the model's own outputs, that can be reported alongside any LLM credit rating so that specification sensitivity is characterized rather than hidden in a single point estimate. Methodologically, the study is pre-registered, uses issuer-clustered resampling and permutation inference for its main uncertainty assessments, and provides a deterministic offline reproducibility capsule (Zenodo, \url{https://doi.org/10.5281/zenodo.21953934}; MIT license) that regenerates every reported number, table, and figure. The study uses only synthetic firm profiles and publicly available SEC filings, involves no human or animal subjects, and discloses its use of generative-AI assistance.

The manuscript is original, is not under consideration elsewhere, and has been approved by both authors, who declare no competing interests and received no external funding. Thank you for your consideration.

\vspace{12pt}
\noindent Sincerely,

\vspace{2pt}
\noindent Robin Chawla\\
Corresponding Author (on behalf of Samir Chincholikar and Robin Chawla)\\
Independent researcher, New York, USA\\
ORCID: \href{https://orcid.org/0009-0007-2807-3948}{0009-0007-2807-3948}\\
\href{mailto:robin.chawla.cse14@iitbhu.ac.in}{robin.chawla.cse14@iitbhu.ac.in}

\end{document}
